\section{Observed cycles and the missing-area identity}\label{sec:completion}
\subsection{Complete images and visible incidence}
\begin{lemma}[Analytic images and unique matching]\label{lem:images}
A nonempty open arc of a compact connected embedded real-analytic
boundary determines its complete image among such boundaries. If $C,D$
are congruent asymmetric compact bodies, exactly one Euclidean isometry
maps $C$ onto $D$.
\end{lemma}
\begin{proof}
For the first assertion consider the points where two candidate curves
have the same germ. This set is open. At a limit point, local analytic
graph representations and the one-variable identity theorem extend
equality of the approaching common arcs, so the set is also closed.
Connectedness gives equality of the complete curves. The convex body
bounded by the curve is then fixed. For the second assertion, the
quotient of any two matching isometries is a symmetry of $C$ and is
therefore the identity.
\end{proof}

The local inverse and Lemma~\ref{lem:images} reconstruct each pair in
its relative placement. Shape separation makes congruence of its
individual members identify precisely their true obstacle types. This
gives the finite visible multigraph. Repeated experiments remain parallel
edges; an edge joining two copies of the same type becomes a loop.
We do not infer identities by comparing the experimental time windows.

\begin{lemma}[Assembly of an arbitrary connected observed component]
\label{lem:assembly}
For any connected component of this multigraph choose a root and a
spanning tree. The records determine one representative of every visible
body in a common plane. Each non-tree edge determines a displacement
$d_e\in\Lambda$. The group $\Gamma=\sum_e\Z d_e$ is independent of
the spanning tree and of all representative and reversal choices, up to
the orthogonal part of one ambient Euclidean isometry.
\end{lemma}
\begin{proof}
Place the root body arbitrarily. For an outward tree edge, match the
source member of its reconstructed pair to the already placed body.
Lemma~\ref{lem:images} gives a unique map; applying that same map to
the partner places the child. In the true periodic table these choices
amount to choosing one lifted representative of each visible type, with
tree labels zero. The labels are used only to justify the construction,
not supplied to it.

For a non-tree edge perform the same source matching. Its target image
$\widehat C_{j,e}$ and the placed target representative $C_j$ are
translates, because their true types agree. This translation is unique;
using the Steiner point $s(C)$ it is
\[
 d_e=s(\widehat C_{j,e})-s(C_j),\qquad
                         \widehat C_{j,e}=C_j+d_e.
\]
It lies in the true period lattice. Reversing a local arclength choice
reflects the whole pair and composes the matching map with that reflection,
leaving its placement unchanged. Reversing an edge changes the sign of
its cycle translation. Changing representative copies adds a vertex
coboundary to the edge translations; these differences cancel on closed
walks. Fundamental cycles generate all closed walks integrally. Thus
the group does not depend on the tree. Parallel repetitions either
repeat a generator or introduce a zero cycle, and cause no problem.
\end{proof}

\subsection{The certificate and its consequences}
\begin{proof}[Proof of Theorem~\ref{thm:completion-main}]
The component identifies exactly a subset $I$ of the true types, because
of shape separation, even though $r$ and the complement of $I$ are
unknown. Lemma~\ref{lem:assembly} gives $\Gamma\subset\Lambda$.
If its rank is two, the index $m=[\Lambda:\Gamma]$ is a positive
integer and $\covol\Gamma=m\covol\Lambda$. Subtracting the
normalization identity \eqref{eq:area-budget} gives
\begin{equation}\label{eq:defect}
 \covol\Gamma-A-S_I=(m-1)\covol\Lambda+S_{I^c},\qquad
 S_{I^c}=\sum_{i\notin I}\area(C_i).
\end{equation}
Both terms on the right are nonnegative. Every obstacle has positive
area. Hence equality to zero is equivalent to $m=1$ and $I^c$ empty.
In that case all representatives have been placed and the full period
group is $\Gamma$, which determines the table. Any other physical
table in the stated class giving these records must have the same
visible images, cycle group and $A$. Applying \eqref{eq:defect} to it
also excludes hidden types and a larger period group. Thus the inverse
is unique even among candidates not initially assumed to be covered.
\end{proof}

\begin{corollary}[A quantitative separation of incomplete data]
\label{cor:defect-gap}
Suppose $\covol\Lambda\ge V_0>0$ and each obstacle orbit has area at
least $a_0>0$. Every connected rank-two observed component has either
$\mathcal D=0$ or
\begin{equation}\label{eq:defect-gap}
                    \mathcal D\ge g_0:=\min(V_0,a_0).
\end{equation}
\end{corollary}
\begin{proof}
If $m>1$, the first term of \eqref{eq:defect} is at least $V_0$.
If some type is unseen, the second term is at least $a_0$.
\end{proof}

\begin{corollary}[Finite positive stopping certificate]\label{cor:stopping}
Any acquisition source returning genuine records may be stopped once
some connected rank-two component has zero defect. No prior claim
that the source retained a specified witness is needed for soundness.
If the accumulated records eventually include a connected covering set
whose cycles generate $\Lambda$, this stopping criterion is eventually
satisfied. Arbitrary deletions can prevent acceptance but cannot turn
an incomplete physical component into an accepted one.
\end{corollary}
\begin{proof}
Soundness is Theorem~\ref{thm:completion-main}. Adding records preserves
a previously covered saturated component and can only join visible
components or add generators. A covering saturated set gives zero
defect when all of its finitely many records are present. Apply the same
theorem to any retained subcollection for the last assertion.
\end{proof}

The converse does not say that every useful list is accepted: a disconnected
list can contain much geometric information but does not register all
its components. The criterion tests components separately. Nor does
nonacceptance identify which acquisition mechanism failed. It simply
withholds a whole-table assertion.

\subsection{Arithmetic without assuming coverage}
For exact displacements pick any independent pair $D=[d_1\ d_2]$.
There is an integer $n_D=|\det D|/\covol\Lambda$ such that every
coordinate vector $z_e=D^{-1}d_e$ belongs to $n_D^{-1}\Z^2$.
This follows because the finite group $\Lambda/D\Z^2$ has order
$n_D$, hence is annihilated by $n_D$. Let $q$ be a common denominator
of the finitely many $z_e$; the least one divides $n_D$. Form the
integer matrix with columns $qz_e$, including $qe_1,qe_2$. If $H$
is a column basis for their subgroup of $\Z^2$, then
\begin{equation}\label{eq:denominator-group}
 \Gamma=\frac1qDH\Z^2,\qquad
 \covol\Gamma=\frac{|\det D|\,|\det H|}{q^2}.
\end{equation}
An elementary Hermite reduction computes $H$. Equivalently its determinant
is the gcd of the two-column minors: factor the integer column matrix
as $HE$, note that the columns of $E$ generate $\Z^2$, and apply
Cauchy--Binet to an integer right inverse of $E$. This gives gcd one
for its minors and therefore gcd $|\det H|$ for the original matrix.
The general algorithms are classical \cite{KB}; the rank-two argument
here is included to fix the column convention.

Crucially, $n_D$ is not computed as $|\det D|/(A+S_I)$: that substitution
would already assume the conclusion that $I$ contains every type.
Section~\ref{sec:uniform} recovers the rational coordinates using only
an upper denominator bound and proves that they can be identified under
noise before the completion test is run.

\subsection{Physical examples with unobserved obstacles}
\begin{proposition}[Both defects occur in analytic tables]\label{prop:examples}
There are analytic asymmetric shape-separated periodic tables and clear
local records realizing each of the following cases: full cycle generation
with an unseen obstacle; complete shape coverage with period index two;
and zero defect without a primitive pair of recorded cycles.
\end{proposition}
\begin{proof}
Work on $\Lambda=\Z^2$ with a sufficiently small asymmetric analytic
body at every lattice point. For example a small multiple of the support
function
\[
 p(\varphi)=1+\epsilon\cos(2\varphi)
                   +\epsilon\cos(3\varphi+\pi/4),\quad
                         0<|\epsilon|<1/11,
\]
is strictly convex since $p+p''>0$. A rotation preserving its nonzero
second and third modes is trivial. A reflection with axis angle $a$
would require both $2a\in\pi\Z$ and $3a+\pi/4\in\pi\Z$, which
is impossible. It is therefore asymmetric.

The shortest bridges associated with the primitive vectors
$(1,0),(1,2),(-1,3)$ are clear for small bodies: the corresponding
center segments contain no other lattice points and have positive
clearance from all other centers. Their pair determinants have absolute
values $2,3,5$, whose gcd is one, so the three cycles generate $\Z^2$
although no pair does. This gives zero defect. Keeping only $(1,0)$
and $(1,2)$ gives index two and defect one, while still observing the
only shape.

Choose a point on the torus away from the finite union of all these
segments and their short experimental tubes, and put a second, differently
scaled asymmetric body there and at all its lattice translates. It can
be made small enough to leave every selected return branch unchanged.
The representatives are noncongruent. Its area $a$ reduces the free area
by exactly $a$. If only the first type is observed using all three
cycles, the cycle group is $\Z^2$ and the defect is $a$, not zero.
With only the first two cycles the defect is $1+a$. These are actual
periodic tables, not only formal area budgets. Strict clearances and
shape inequalities persist under small suitable perturbations.
\end{proof}
